\section{Preliminary Results}
\label{chap:proposal:preliminary}

The first semester was dedicated to establishing a robust foundation for the dissertation, conducting the State of the Art, while also performing some  preliminary work and exploratory investigations. These initial activities were essential for identifying certain limitations and establishing the groundwork for the functionalities to be implemented in the second semester. This section focuses on the preliminary activities undertaken during the first semester.

The initial phase of this preliminary work involved the deployment and integration of OpenCAPIF for securing a CAMARA \acs{QoD} \acs{API} stub, generated from its Open\acs{API} spec. This contributed to the familiarization with the CAMARA initiative and the \acs{QoD} \acs{API}, marking an important milestone in the development process with the integration of OpenCAPIF, as integrating OpenCAPIF with other \acsp{API} will follow the same approach.

The next task involved the development of the Slice Manager stub. This proved to be a particularly demanding task, as the author did not have access to the source code due to its proprietary nature, requiring the implementation of the \acs{API} solely from a limited set of example interactions (request and response). This extended the development time and  several additional meetings to clarify requirements.

Afterwards, the project referenced in Section \ref{chap:proposal:implementation}, which integrated \acs{GSMA} Open Gateway \acsp{API} and the NEFSim, was introduced with the goal of integrating its slicing functionalities with the Slice Manager stub. However, some inconsistencies regarding the integration were found, which required the author to focus on the additional task of validating and resolving these interactions. This effort is still under progress.

In paralell, it was necessary to become familiar with \acs{ITAv}'s infrastructure, specifically its \acs{OSS}: \acl{OSL}. This involded understanding the \acs{TMF}'s concepts of service and resource, deploying an instance of \acs{OSL} and exposing a Helm chart in \acs{OSL}'s Service Catalog, allowing the automated provisioning of instances upon request. Subsequently, with the goal of managing the lifecycle of slices and \acsp{UE} via the Slice Manager (currently stub), the generic controller capabilities of \acs{OSL} were explored. This culminated in the development of a Slice Generic Controller and a lifecycle management handler responsible for the creation phase of a slice by interacting with the Slice Manager, when provisioned through \acs{OSL}~\cite{OSLSlicingGcontrollerGithub}. This milestone is particularly significant, as the same approach will be applied to update and deletion operations, not only for slices but also for \acsp{UE}, which are included among the first tasks outlined in the implementation plan for the second semester in Section \ref{chap:proposal:implementation}.
